% ja-latex.tex — 和文の LaTeX (lualatex) の PDF の共通の前置き (2026-11-18．yaml の include-in-header から読む)
%   lecture / report-latex / handout-latex / note の pdf が使う．前は各 yaml に同じ文を書いていた．
\usepackage{luatexja-preset}
% luatexja-fontspec (TeX Live 2026，2026-09 以降) が regexpatch を読み，後で xpatch を読むと \xpatchcmd が二重定義の誤りになる．
% regexpatch は xpatch の命令も持つので，xpatch は読んだことにする (amsthm を読む文書で起きた．NOTES.org)
\makeatletter\@ifpackageloaded{regexpatch}{\@ifpackageloaded{xpatch}{}{\expandafter\def\csname ver@xpatch.sty\endcsname{2020/03/25 v0.3a (provided by regexpatch)}}}{}\makeatother
\usepackage{mathtools}
% 式の番号は名前 (eq-) を付けた式と \eqref した式だけ (html と同じ．ox-quarto-ext-math.el)
\mathtoolsset{showonlyrefs,showmanualtags}
